\documentclass{article}

\usepackage{fontspec}
\usepackage{polyglossia}
\setdefaultlanguage{greek}

\usepackage{geometry}
\geometry{a4paper, margin=1in}

\usepackage{amssymb, amsmath}

\usepackage{enumitem}
\setlist[enumerate]{label={$\square$ \arabic*.}}

\setmainfont{Liberation Serif}
\newfontfamily\greekfont{Liberation Serif}

\begin{document}
\section[1]{Κρούσεις και Ταλαντώσεις}

\subsection{Ιδανικά Ελατήρια}

\begin{enumerate}

      \item
            Τι αναφέρει η Αρχή Διατήρησης της Ενέργειας (Α.Δ.Ε.); Αν στο σώμα ασκούνται μόνο συντηρητικές δυνάμεις ή ασκούνται και μη συντηρητικές δυνάμεις που έχουν μηδενικό έργο, σε ποια βασική αρχή καταλήγει η Α.Δ.Ε.; (σελ. 15)

      \item
            Πότε ορίζεται η έννοια της θέσης φυσικού μήκους (Θ.Φ.Μ.) για ένα ελατήριο; (σελ. 15)

      \item
            Αν όλες οι υπόλοιπες δυνάμεις που ασκούνται στο σώμα είναι σταθερές και η μόνη δύναμη που μεταβάλλεται είναι η δύναμη του ελατηρίου, μπορούμε να χρησιμοποιήσουμε τις χρονικές εξισώσεις της ευθύγραμμης ομαλά μεταβαλλόμενης κίνησης για να μελετήσουμε την κίνηση του σώματος; Να αιτιολογήσετε την απάντησή σας. (σελ. 15)

      \item
            Με ποιον τύπο μπορούμε να υπολογίσουμε την τριβή ολίσθησης; Με ποιον τρόπο υπολογίζουμε την κάθετη δύναμη στήριξης που δέχεται ένα σώμα από το επίπεδο πάνω στο οποίο κινείται; (σελ. 17)

      \item
            Με ποιον τύπο υπολογίζουμε το έργο της τριβής ολίσθησης, όταν η τριβή είναι σταθερή; Με τι ισούται η θερμότητα που εκλύεται εξαιτίας της τριβής ολίσθησης; (σελ. 17)

      \item
            Ο 2ος νόμος του Νεύτωνα μπορεί να εφαρμοστεί όταν η συνισταμένη δύναμη μεταβάλλεται; Τι προκύπτει για την επιτάχυνση του σώματος όταν η συνισταμένη δύναμη που δέχεται το σώμα μεταβάλλεται; (σελ. 17)

      \item
            Τι ονομάζουμε θέση ισορροπίας ενός σώματος; Πώς μπορούμε να βρούμε την απόσταση της θέσης ισορροπίας από τη θέση φυσικού μήκους του ελατηρίου; (σελ. 20)

      \item
            Κατά τη μελέτη της κίνησης ενός συστήματος σώμα - ελατήριο, πότε κατά τον υπολογισμό της μηχανικής ενέργειας του συστήματος πρέπει να συμπεριλάβουμε και τη δυναμική ενέργεια του σώματος λόγω βαρύτητας; (σελ. 20)

      \item
            Πώς υπολογίζουμε τη δυναμική ενέργεια λόγω βαρύτητας ενός σώματος; (σελ. 20)

      \item
            Ποια είναι η σχέση υπολογισμού του ρυθμού μεταβολής της κινητικής ενέργειας ενός υλικού σημείου; Ποια είναι η απόδειξη της σχέσης αυτής; (σελ. 22-23)

      \item
            Ποια πληροφορία μας δίνει η αλγεβρική τιμή ενός διανύσματος; Πότε ένας διανυσματικός τύπος μπορεί να μετατραπεί σε αλγεβρικό; (σελ. 23)

      \item
            Με ποιον τύπο υπολογίζουμε το έργο μιας σταθερής δύναμης που το σημείο εφαρμογής της μετατοπίζεται ευθύγραμμα; (σελ. 26)

      \item
            Με ποιον τύπο μπορούμε να υπολογίζουμε το ρυθμό μεταβολής της δυναμικής ενέργειας του ιδανικού ελατηρίου; Ποια είναι η απόδειξη του τύπου αυτού; (σελ. 26)

      \item
            Με ποιον τύπο μπορούμε να υπολογίσουμε σε κάθε περίπτωση το έργο του βάρους ενός σώματος σταθερής μάζας; Να εξηγήσετε αναλυτικά τι σημαίνει το κάθε μέγεθος στον τύπο που γράψατε. Πότε το έργο του βάρους ενός σώματος είναι θετικό και πότε αρνητικό; (σελ. 29)

      \item
            Ποιος είναι ο ενεργειακός ρόλος της δύναμης του ελατηρίου και της βαρυτικής δύναμης; Τι επίδραση έχουν στη μηχανική ενέργεια; Με τι ισούται το έργο των δυνάμεων αυτών σε μια κλειστή διαδρομή; (σελ. 31)
\end{enumerate}

\subsection{Θεωρία Κρούσεων}

\begin{enumerate}
      \item Τι ονομάζουμε κρούση στη μηχανική; (σελ. 47)

      \item Τι ονομάζουμε κρούση στην ατομική και πυρηνική φυσική; (σελ. 47)

      \item Ποια κρούση ονομάζεται κεντρική, ποια έκκεντρη και ποια πλάγια; (σελ. 47 - 48)

      \item Πότε μια κρούση ονομάζεται ελαστική και πότε ανελαστική; (σελ. 49)

      \item Πότε μια κρούση ονομάζεται πλαστική και τι συμβαίνει με την ενέργεια του συστήματος; (σελ. 49)

      \item Δύο σφαίρες με μάζες $m_1$ και $m_2$ κινούνται πάνω σε λείο οριζόντιο δάπεδο και συγκρούονται κεντρικά και ελαστικά έχοντας ελάχιστα πριν τη σύγκρουση ταχύτητες αλγεβρικής τιμής $v_1$ και $v_2$ αντίστοιχα. Χρησιμοποιώντας τις αρχές διατήρησης που ισχύουν στην κρούση αυτή να γράψετε και στη συνέχεια να αποδείξετε τις σχέσεις υπολογισμού των αλγεβρικών τιμών των ταχυτήτων των δύο σφαιρών αμέσως μετά την κρούση. (σελ. 49 - 50)

      \item Δύο σφαίρες που έχουν ίσες μάζες συγκρούονται κεντρικά και ελαστικά. Χρησιμοποιώντας τις αρχές διατήρησης που ισχύουν σε μια τέτοια κρούση να αποδείξετε ότι οι δύο σφαίρες ανταλλάσσουν ταχύτητες. (σελ. 51)

      \item Μια σφαίρα μάζας $m_1$ συγκρούεται κεντρικά και ελαστικά με μια ακίνητη σφαίρα μάζας $m_2$ για την οποία ισχύει ότι $m_2 \gg m_1$. Να αποδείξετε ότι η σφαίρα μάζας $m_1$ θα αποκτήσει, εξαιτίας της ελαστικής κρούσης, αντίθετη ταχύτητα από αυτή που είχε ελάχιστα πριν την κρούση. (σελ. 53)

      \item Μια σφαίρα μάζας $m_1$ κινείται με ταχύτητα μέτρου $v_1$ και συγκρούεται κεντρικά και ελαστικά με μια ακίνητη σφαίρα μάζας $m_2$, για την οποία ισχύει ότι $m_2 \ll m_1$. Να αποδείξετε ότι η αρχικά ακίνητη σφαίρα μάζας $m_2$ θα αποκτήσει, εξαιτίας της ελαστικής κρούσης, ταχύτητα μέτρου $2v_1$. (σελ. 53 - 54)

      \item Μια σφαίρα συγκρούεται πλάγια και ελαστικά μ' έναν κατακόρυφο τοίχο. Να αποδείξετε ότι η γωνία πρόσπτωσης της σφαίρας στον τοίχο ισούται με τη γωνία ανάκλασης της από αυτόν. (σελ. 54 - 55)

      \item Με ποιον τύπο μπορούμε να υπολογίσουμε την απώλεια κινητικής ενέργειας σε μια ανελαστική κρούση δύο σωμάτων; (σελ. 55)
\end{enumerate}

\subsection{Κεντρικές Μετωπικές Κρούσεις}

\begin{enumerate}

      \item
            Ποιοι είναι οι τύποι με τους οποίους υπολογίζουμε τις ταχύτητες των σωμάτων αμέσως μετά την κεντρική και ελαστική κρούση δύο σφαιρών; Τι πρέπει να προσέχεις όταν αντικαθιστάς τις τιμές των μεγεθών; (σελ. 91)

      \item
            Με ποιον τύπο μπορούμε να υπολογίσουμε τη μεταβολή της ορμής ενός σώματος εξαιτίας της κεντρικής κρούσης στην οποία εμπλέκεται; Γιατί τον τύπο αυτόν μπορούμε να τον γράψουμε αλγεβρικά; (σελ. 91)

      \item
            Όταν η ορμή ενός συστήματος δύο σωμάτων παραμένει σταθερή, ποια είναι η σχέση που συνδέει τη μεταβολή της ορμής του ενός σώματος με τη μεταβολή της ορμής του άλλου; Μπορείς να αποδείξεις τη σχέση αυτή; (σελ. 91-92)

      \item
            Σε μια ελαστική κρούση μεταξύ δύο σωμάτων, ποια είναι η σχέση που συνδέει τη μεταβολή της κινητικής ενέργειας του ενός σώματος με τη μεταβολή της κινητικής ενέργειας του άλλου; Μπορείς να αποδείξεις τη σχέση αυτή; (σελ. 94)

      \item
            Με ποιον τύπο μπορούμε να υπολογίσουμε το ποσοστό επί τοις εκατό της μεταβολής της κινητικής ενέργειας ενός σώματος εξαιτίας μιας κρούσης; (σελ. 94)

      \item
            Σε μια ελαστική κρούση μεταξύ δύο σωμάτων, με ποιον τύπο μπορούμε να υπολογίσουμε το ποσοστό επί τοις εκατό της κινητικής ενέργειας ενός σώματος που μεταβιβάζεται στο άλλο σώμα; (σελ. 96)

      \item
            Με ποιον τύπο υπολογίζουμε την απώλεια της μηχανικής ενέργειας ενός συστήματος σωμάτων που συγκρούονται ανελαστικά; (σελ. 99)

      \item
            Σε μια ανελαστική κρούση μεταξύ δύο σωμάτων, με ποιον τύπο μπορούμε να υπολογίσουμε το ποσοστό επί τοις εκατό της αρχικής κινητικής ενέργειας του συστήματος που χάθηκε από το σύστημα εξαιτίας της κρούσης; (σελ. 99)

      \item
            Τι ελέγχουμε αν θέλουμε να διερευνήσουμε αν μια κρούση μεταξύ δύο σωμάτων είναι ελαστική ή ανελαστική; (σελ. 99)

      \item
            Πώς χαρακτηρίζεται η κρούση (ανελαστική ή ελαστική) όταν ένα βλήμα εισέρχεται σε ένα σώμα και αμέσως μετά την κρούση εξέρχεται από αυτό; (σελ. 102)

      \item
            Ποια είναι τα δύο βασικά εργαλεία που χρησιμοποιούμε για τη μελέτη μιας οποιασδήποτε ανελαστικής κεντρικής (ή μετωπικής) κρούσης; Με ποιον τρόπο υπολογίζεται η μεταβολή της ορμής για κάθε σώμα που εμπλέκεται σε μια ανελαστική κεντρική (ή μετωπική) κρούση; (σελ. 102)

      \item
            Τι ανταλλάσσουν δύο σώματα που έχουν ίσες μάζες όταν τα σώματα αυτά συγκρούονται κεντρικά και ελαστικά; (σελ. 105)

      \item
            Πότε είναι προτιμότερο (πιο εύκολο) να χρησιμοποιούμε ενεργειακό εργαλείο για τη μελέτη μιας ευθύγραμμης ομαλά μεταβαλλόμενης κίνησης που εκτελεί ένα σώμα πριν ή μετά από μια κρούση; (σελ. 105)

      \item
            Όταν θέλεις να μελετήσεις την κίνηση που εκτελεί ένα σώμα μετά από μια κρούση, τι πρέπει να προσέχεις όταν φτιάχνεις το σχήμα σου; Ποια είναι η αρχική ταχύτητα του σώματος της κίνησης αυτής; (σελ. 105)

      \item
            Με ποιον τύπο μπορείς να υπολογίσεις το έργο της δύναμης του ελατηρίου; (σελ. 107)

      \item
            Με τι ισούται ο ρυθμός μεταβολής της ορμής ενός σώματος; Τι πρέπει να προσέξεις για το ρυθμό αυτό, όταν το σώμα κινείται ευθύγραμμα; (σελ. 111)

      \item
            Με ποιον τύπο μπορείς να υπολογίσεις το ρυθμό μεταβολής της κινητικής ενέργειας ενός υλικού σημείου; Μπορείς να αποδείξεις τον τύπο αυτό; (σελ. 112)

      \item
            Τι πρέπει να προσέχεις όταν εφαρμόζεις την Α.Δ.Ε. ή την Α.Δ.Μ.Ε. στην περίπτωση που εμπλέκεται στην κίνηση του σώματος ιδανικό ελατήριο το οποίο είναι κατακόρυφο ή βρίσκεται σε κεκλιμένο επίπεδο; (σελ. 116)

      \item
            Μπορεί ένα σώμα που έχει μηδενική ταχύτητα μια χρονική στιγμή να έχει μη μηδενική επιτάχυνση την ίδια χρονική στιγμή; Εξαρτάται η επιτάχυνση ενός σώματος σε μια χρονική στιγμή από την ταχύτητα που έχει το σώμα την ίδια χρονική στιγμή; (σελ. 116)

      \item
            Με τι ισούται το έργο της τάσης του νήματος σε μια κυκλική διαδρομή όπου η τάση βρίσκεται συνεχώς στη διεύθυνση της ακτίνας; Μπορείς να αιτιολογήσεις την απάντησή σου; (σελ. 122)

      \item
            Με ποιον τύπο μπορούμε να υπολογίσουμε σε κάθε περίπτωση κίνησης το έργο του βάρους ενός σώματος; Να ερμηνεύσεις τα μεγέθη του τύπου αυτού και να εξηγήσεις πώς επιλέγουμε το πρόσημο του τύπου. (σελ. 124)

      \item
            Με ποιον τρόπο μπορούμε να σχεδιάσουμε και να υπολογίσουμε την κεντρομόλο δύναμη που δέχεται ένα σώμα σε ένα σημείο της κυκλικής κίνησης που εκτελεί; (σελ. 128)

      \item
            Ποιες σχέσεις περιλαμβάνει ο 2ος νόμος του Νεύτωνα στην περίπτωση της κυκλικής κίνησης; (σελ. 132)

      \item
            Τι σημαίνει ότι ένα σώμα που κινείται κυκλικά σε κατακόρυφο επίπεδο εκτελεί ανακύκλωση; Με ποιον τρόπο μπορούμε να υπολογίσουμε την ελάχιστη ταχύτητα που πρέπει να έχει ένα σώμα στο ανώτατο σημείο της τροχιάς του, ώστε να εκτελέσει ανακύκλωση; (σελ. 132)

      \item
            Με ποιον τύπο μπορούμε να υπολογίσουμε το ρυθμό μεταβολής της κινητικής ενέργειας ενός υλικού σημείου σε ένα σημείο της κυκλικής κίνησης που εκτελεί; Μπορείς να αποδείξεις τον τύπο αυτό; (σελ. 133)

      \item
            Ποια είναι τα χαρακτηριστικά του συστήματος δύο σωμάτων που είναι ελεύθερα να κινηθούν σε λείο οριζόντιο δάπεδο και είναι συνδεδεμένα με οριζόντιο ιδανικό ελατήριο; (σελ. 137-138)
\end{enumerate}

\subsection{Πλάγιες Κρούσεις}

\begin{enumerate}
      \item
            Όταν ένα μικρό σώμα συγκρούεται πλάγια και ελαστικά με λείο τοίχο ή δάπεδο, ποια είναι η σχέση της γωνίας πρόσπτωσης με τη γωνία ανάκλασης; Μπορείς το συμπέρασμα αυτό να το χρησιμοποιήσεις χωρίς απόδειξη; (σελ. 174)

      \item
            Όταν ένα σώμα συγκρούεται με λείο τοίχο ή δάπεδο, ποια είναι η διεύθυνση της δύναμης που δέχεται το σώμα αυτό από τον τοίχο ή το δάπεδο; Τι σημαίνει μέση δύναμη στην περίπτωση αυτή; (σελ. 174)

      \item
            Όταν δύο σώματα συγκρούονται πλάγια (π.χ. δύο σφαίρες που κινούνται σε λείο οριζόντιο δάπεδο), ποια είναι η διαφορά στην εφαρμογή της Α.Δ.Ο. σε σχέση με την περίπτωση που τα δύο αυτά σώματα συγκρούονται κεντρικά; (σελ. 178)

      \item
            Μπορείς να περιγράψεις πώς εφαρμόζεται η Α.Δ.Ο. με ανάλυση σε άξονες σε μια πλάγια κρούση μεταξύ δύο σφαιρών; (σελ. 178 - 179)

      \item
            Μπορείς να περιγράψεις έναν χρήσιμο τρόπο εφαρμογής της Α.Δ.Ο. σε μια πλάγια κρούση μεταξύ δύο σφαιρών όταν τα διανύσματα που "συμμετέχουν" στην Α.Δ.Ο. είναι τρία; (σελ. 179)

      \item
            Υπάρχει κάποια διαφορά στο πώς αντιμετωπίζουμε την ενέργεια στις πλάγιες κρούσεις σε σχέση με το πώς την αντιμετωπίζουμε στις κεντρικές κρούσεις; (σελ. 181)

      \item
            Πώς αντιμετωπίζουμε την περίπτωση πλάγιας κρούσης στην οποία το σύστημα στη διεύθυνση του ενός άξονα δεν είναι μονωμένο; Μπορείς να αναγνωρίσεις πότε συμβαίνει αυτό; (σελ. 183)
\end{enumerate}

\subsection{Θεωρία Μηχανικές Ταλαντώσεις}

\begin{enumerate}
      \item Ποια φαινόμενα ονομάζονται περιοδικά; (σελ. 201)

      \item Τι ονομάζεται περίοδος και τι συχνότητα ενός περιοδικού φαινομένου και ποια είναι η μαθηματική σχέση που συνδέει τα δύο αυτά μεγέθη; (σελ. 201)

      \item Ποιο μέγεθος ονομάζεται γωνιακή συχνότητα και ποια είναι η μονάδα μέτρησής της; (σελ. 201)

      \item Τι ονομάζεται πλάτος και τι θέση ισορροπίας μιας ταλάντωσης; (σελ. 202)

      \item Αν ένα σώμα εκτελεί απλή αρμονική ταλάντωση και η χρονική εξίσωση της απομάκρυνσης από τη θέση ισορροπίας του είναι της μορφής $x = A \eta\mu \omega t$, τότε ποιες είναι οι χρονικές εξισώσεις της ταχύτητας και της επιτάχυνσής του; (σελ. 203-204)

      \item Να γράψετε και να αποδείξετε την αλγεβρική σχέση που συνδέει την ταχύτητα $v$ και την απομάκρυνση $x$ από τη θέση ισορροπίας για ένα σώμα που εκτελεί απλή αρμονική ταλάντωση. (σελ. 204)

      \item Να γράψετε και να αποδείξετε την αλγεβρική σχέση που συνδέει την επιτάχυνση $a$ και την απομάκρυνση $x$ από τη θέση ισορροπίας για ένα σώμα που εκτελεί απλή αρμονική ταλάντωση. (σελ. 205)

      \item Να γράψετε την αλγεβρική σχέση που συνδέει τη συνισταμένη δύναμη $\Sigma F$ με την απομάκρυνση $x$ από τη θέση ισορροπίας για ένα σώμα το οποίο εκτελεί απλή αρμονική ταλάντωση και στη συνέχεια να σχεδιάσετε ποιοτικά τη γραφική της παράσταση. (σελ. 207)

      \item Με τι ισούται η σταθερά επαναφοράς και από τι εξαρτάται; Ποια είναι η μονάδα μέτρησής της στο σύστημα μονάδων S.I.; (σελ. 207)

      \item Γιατί η συνισταμένη δύναμη που δέχεται ένα σώμα το οποίο εκτελεί α.α.τ. ονομάζεται και δύναμη επαναφοράς; (σελ. 207)

      \item Ποια σχέση συνδέει την περίοδο μιας απλής αρμονικής ταλάντωσης και τη σταθερά επαναφοράς; (σελ. 207)

      \item Με το τετράγωνο ποιου μεγέθους είναι ανάλογη η ενέργεια μιας απλής αρμονικής ταλάντωσης; (σελ. 208)

      \item Ποια είναι η σχέση υπολογισμού της δυναμικής ενέργειας μιας απλής αρμονικής ταλάντωσης; Πώς αποδεικνύεται η σχέση αυτή; (σελ. 208)

      \item Αν η χρονική εξίσωση της απομάκρυνσης από τη θέση ισορροπίας για ένα σώμα που εκτελεί α.α.τ. είναι της μορφής $x = A \eta\mu \omega t$, να γράψετε τις αντίστοιχες χρονικές εξισώσεις της κινητικής και της δυναμικής ενέργειας και να σχεδιάσετε ποιοτικά τις γραφικές τους παραστάσεις για τη χρονική διάρκεια από $0$ έως $T$ (περίοδος) σε κοινό σύστημα αξόνων. (σελ. 209-210)

      \item Να γράψετε τις εξισώσεις $U = f(x)$ και $K = f(x)$ για μια α.α.τ. και στη συνέχεια να σχεδιάσετε ποιοτικά τις αντίστοιχες γραφικές παραστάσεις σε κοινό σύστημα αξόνων. (σελ. 210)

      \item Να γράψετε τις εξισώσεις $U = f(v)$ και $K = f(v)$ για μια α.α.τ. και στη συνέχεια να σχεδιάσετε ποιοτικά τις αντίστοιχες γραφικές παραστάσεις σε κοινό σύστημα αξόνων. (σελ. 211)

      \item Χρησιμοποιώντας την Α.Δ.Ε. για την ταλάντωση που εκτελεί ένα σώμα να αποδείξετε την αλγεβρική σχέση που συνδέει την ταχύτητα $v$ και την απομάκρυνση $x$ από τη θέση ισορροπίας του. (σελ. 211)

      \item Να γράψετε τις χρονικές εξισώσεις $x = f(t)$, $v = f(t)$ και $a = f(t)$ για μια απλή αρμονική ταλάντωση που έχει αρχική φάση. (σελ. 211)

      \item Τι ονομάζεται φάση της ταλάντωσης; Ποια είναι η γραφική παράσταση της φάσης της ταλάντωσης σε συνάρτηση με το χρόνο; (σελ. 212)

      \item Πότε η ταλάντωση ενός σώματος δεν έχει αρχική φάση; (σελ. 212)

      \item Τι σημαίνει ότι δύο αρμονικά εναλλασσόμενα μεγέθη έχουν διαφορά φάσης; (σελ. 212)

      \item Πώς μπορούμε να μελετήσουμε τη χρονική εξίσωση $x = f(t)$ μιας απλής αρμονικής ταλάντωσης με τη βοήθεια ενός περιστρεφόμενου διανύσματος; (σελ. 213)

      \item Πώς μπορούμε να βρούμε την αρχική φάση μιας απλής αρμονικής ταλάντωσης με τη βοήθεια του περιστρεφόμενου διανύσματος; (σελ. 214)

      \item Πώς μπορεί να παρασταθεί η διαφορά φάσης μεταξύ δύο αρμονικά εναλλασσόμενων μεγεθών χρησιμοποιώντας τα αντίστοιχα περιστρεφόμενα διανύσματά τους; (σελ. 214)

      \item Τι ονομάζεται θέση φυσικού μήκους (Θ.Φ.Μ.) ενός ελατηρίου; (σελ. 215)

      \item Πότε ασκεί δύναμη το ελατήριο και με ποιον τύπο υπολογίζουμε το μέτρο της; (σελ. 215)

      \item Με ποιον τύπο μπορούμε να υπολογίσουμε τη δυναμική ενέργεια λόγω παραμόρφωσης ενός ελατηρίου; Να περιγράψετε μεγέθη στον τύπο που γράψατε. (σελ. 215)

      \item Με ποιον τύπο μπορούμε να υπολογίσουμε το έργο της δύναμης ενός ελατηρίου; Είναι η δύναμη του ελατηρίου μια συντηρητική δύναμη; (σελ. 216)

      \item Να αποδείξετε ότι αν ένα σώμα είναι δεμένο στο ένα άκρο ιδανικού ελατηρίου σταθεράς $k$ και κινείται πάνω σε λείο οριζόντιο δάπεδο με τη δράση του βάρους του, της δύναμης από το δάπεδο και της δύναμης του ελατηρίου, τότε εκτελεί μια απλή αρμονική ταλάντωση με σταθερά επαναφοράς $D = k$. (σελ. 216-217)

      \item Να αποδείξετε ότι αν ένα σώμα είναι δεμένο στο ένα άκρο κατακόρυφου ιδανικού ελατηρίου σταθεράς $k$ και κινείται με τη δράση του βάρους του και της δύναμης του ελατηρίου, τότε εκτελεί μια απλή αρμονική ταλάντωση με σταθερά επαναφοράς $D = k$. (σελ. 217)

      \item Να αποδείξετε ότι αν ένα σώμα είναι δεμένο στο ένα άκρο ιδανικού ελατηρίου σταθεράς $k$ και κινείται πάνω σε λείο κεκλιμένο επίπεδο με τη δράση του βάρους του, της δύναμης από το επίπεδο και της δύναμης του ελατηρίου, τότε εκτελεί μια απλή αρμονική ταλάντωση με σταθερά επαναφοράς $D = k$. (σελ. 217-218)

      \item Με ποιον τύπο υπολογίζεται η περίοδος της απλής αρμονικής ταλάντωσης ενός συστήματος μάζα - ελατήριο; (σελ. 218)

      \item Ένα σώμα εκτελεί φθίνουσα ταλάντωση και η δύναμη αντίστασης στην κίνησή του είναι της μορφής $F' = -bv$. Πώς ονομάζεται η σταθερά $b$ και από τι εξαρτάται; Ποια είναι η μονάδα μέτρησής της σταθεράς $b$ στο σύστημα μονάδων S.I.; (σελ. 219)

      \item Αν αυξήσουμε την πίεση σ' ένα δοχείο με αέρα, τότε θα αυξηθεί ή θα μειωθεί η σταθερά $b$; (σελ. 219)

      \item Ποια σχέση συνδέει την επιτάχυνση, την ταχύτητα και την απομάκρυνση μια δεδομένη χρονική στιγμή στη φθίνουσα ταλάντωση; Πώς αποδεικνύεται η σχέση αυτή; (σελ. 219)

      \item Να σχεδιάσετε τη γραφική παράσταση της απομάκρυνσης $x$ από τη θέση ισορροπίας ενός σώματος που εκτελεί φθίνουσα ταλάντωση σε συνάρτηση με το χρόνο για δύο μικρές τιμές $b_1$ και $b_2$ της σταθεράς απόσβεσης, με $b_2 > b_1$. (σελ. 220)

      \item Να γράψετε τη χρονική εξίσωση του πλάτους $A$ για ένα σώμα που εκτελεί φθίνουσα ταλάντωση και δέχεται δύναμη αντίστασης στην κίνησή του της μορφής $F' = -bv$ και στη συνέχεια να σχεδιάσετε τη γραφική παράσταση της εξίσωσης αυτής για δύο μικρές διαφορετικές τιμές $b_1$ και $b_2$ της σταθεράς απόσβεσης, με $b_2 > b_1$. (σελ. 220)

      \item Με ποιον τύπο μπορούμε να υπολογίσουμε τον χρόνο ημιζωής (ή ημίσειας ζωής) του πλάτους $A$ για ένα σώμα που εκτελεί φθίνουσα ταλάντωση και δέχεται δύναμη αντίστασης στην κίνησή του της μορφής $F' = -bv$; Να αποδείξετε τον τύπο αυτό. (σελ. 220-221)

      \item Για ένα σώμα που εκτελεί φθίνουσα ταλάντωση και δέχεται δύναμη αντίστασης στην κίνησή του της μορφής $F' = -bv$, να αποδείξετε ότι το πηλίκο δύο διαδοχικών μέγιστων απομακρύνσεων προς την ίδια κατεύθυνση παραμένει σταθερό. (σελ. 221-222)

      \item Ποια είναι η σχέση της περιόδου $T$ της φθίνουσας ταλάντωσης που εκτελεί ένα σύστημα και της ιδιοπεριόδου $T_0$ του συστήματος αυτού; Πότε οι δύο αυτές περίοδοι μπορεί να θεωρηθεί ότι ταυτίζονται; (σελ. 222)

      \item Κατά τη διάρκεια μιας φθίνουσας ταλάντωσης, όπου το πλάτος μειώνεται εκθετικά με το χρόνο, η συχνότητα και η περίοδος της φθίνουσας ταλάντωσης παραμένουν σταθερές ή ελαττώνονται και αυτές με την πάροδο του χρόνου; (σελ. 222)

      \item Η ενέργεια της φθίνουσας ταλάντωσης με το άθροισμα ποιων δύο επιμέρους ενεργειών είναι ίση; Η ενέργεια της φθίνουσας ταλάντωσης είναι σταθερή ή ελαττώνεται με την πάροδο του χρόνου; (σελ. 222)

      \item Στην περίπτωση μικρής τιμής της σταθεράς απόσβεσης με ποιον τύπο μπορούμε να υπολογίσουμε προσεγγιστικά την ενέργεια της φθίνουσας ταλάντωσης; Να αποδείξετε τον τύπο αυτό. (σελ. 222)

      \item Πώς ονομάζεται η συχνότητα μιας ελεύθερης και αμείωτης ταλάντωσης; Με ποιον τύπο την υπολογίζουμε; (σελ. 223)

      \item Ποια ταλάντωση ονομάζεται εξαναγκασμένη; (σελ. 224)

      \item Με τι ισούται η συχνότητα και η περίοδος της εξαναγκασμένης ταλάντωσης που εκτελεί ένα σύστημα; (σελ. 224)

      \item Αν η συχνότητα του διεγέρτη παραμένει σταθερή, τότε το πλάτος της εξαναγκασμένης ταλάντωσης παραμένει σταθερό ή αλλάζει με την πάροδο του χρόνου; (σελ. 224)

      \item Να σχεδιάσετε τη γραφική παράσταση με την οποία φαίνεται πώς μεταβάλλεται το πλάτος μιας εξαναγκασμένης ταλάντωσης σε συνάρτηση με τη συχνότητα του διεγέρτη για συγκεκριμένη τιμή της σταθεράς απόσβεσης. (σελ. 225)

      \item Αν η σταθερά απόσβεσης είναι μικρή, τότε για ποια συχνότητα του διεγέρτη επιτυγχάνεται η μεγιστοποίηση του πλάτους; (σελ. 225)

      \item Αν ένα σύστημα εκτελεί εξαναγκασμένη ταλάντωση και αυξήσουμε τη σταθερά απόσβεσης χωρίς να αλλάξουμε τη συχνότητα του διεγέρτη, τότε το πλάτος της εξαναγκασμένης ταλάντωσης θα αυξηθεί, θα μειωθεί ή θα παραμεινει σταθερό; (σελ. 225)

      \item Να σχεδιάσετε τη γραφική παράσταση με την οποία φαίνεται πώς μεταβάλλεται το πλάτος μιας εξαναγκασμένης ταλάντωσης σε συνάρτηση με τη συχνότητα του διεγέρτη για διαφορετικές τιμές της σταθεράς απόσβεσης. (σελ. 225)

      \item Ποια συμπεράσματα προκύπτουν για την ενέργεια ενός συστήματος που εκτελεί εξαναγκασμένη ταλάντωση; (σελ. 226)

      \item Πότε λέμε ότι ένα σύστημα, το οποίο εκτελεί εξαναγκασμένη ταλάντωση, βρίσκεται σε κατάσταση συντονισμού; Τι συμβαίνει στην κατάσταση αυτή; (σελ. 226)
\end{enumerate}

\subsection{Απλή Αρμονική Ταλάντωση}

\begin{enumerate}
      \item Με ποιον τύπο υπολογίζεται η μέγιστη ταχύτητα και η μέγιστη επιτάχυνση ενός σώματος που εκτελεί απλή αρμονική ταλάντωση; Σε ποια θέση επιτυγχάνεται κατά μέτρο η μέγιστη ταχύτητα και η μέγιστη επιτάχυνση; (σελ. 259)

      \item Ένα σώμα που εκτελεί απλή αρμονική ταλάντωση διέρχεται κάποια χρονική στιγμή από ένα σημείο Α της ταλάντωσης του. Από ποιο σημείο θα διέλθει το σώμα αυτό και με ποια ταχύτητα τη χρονική στιγμή που ολοκληρώνει μία πλήρη ταλάντωση; (σελ. 259)

      \item Πόσο είναι το διάστημα που διανύει ένα σώμα στη χρονική διάρκεια μίας πλήρους ταλάντωσης; Εξαρτάται το διάστημα αυτό από ποιο σημείο ξεκινά η ταλάντωση; (σελ. 259)

      \item Ποιες είναι οι χρονικές εξισώσεις της απομάκρυνσης, της ταχύτητας και της επιτάχυνσης για ένα σώμα που εκτελεί απλή αρμονική ταλάντωση αν τη χρονική στιγμή $t = 0$ το σώμα διέρχεται από τη θέση ισορροπίας του με θετική ταχύτητα; (σελ. 259 - 260)

      \item Ποια είναι η σχέση που συνδέει την επιτάχυνση ενός σώματος που εκτελεί απλή αρμονική ταλάντωση με την απομάκρυνσή του από τη θέση ισορροπίας του; Εξαρτάται η εξίσωση αυτή από τη θέση που βρίσκεται το σώμα τη χρονική στιγμή $t = 0$; (σελ. 263)

      \item Ποια είναι η σχέση που συνδέει την ταχύτητα ενός σώματος που εκτελεί απλή αρμονική ταλάντωση με την απομάκρυνσή του από τη θέση ισορροπίας του; Με ποιον τρόπο μπορούμε να αποδείξουμε τη σχέση αυτή; (σελ. 263)

      \item Ποια είναι η σχέση που συνδέει την επιτάχυνση ενός σώματος που εκτελεί απλή αρμονική ταλάντωση με την ταχύτητά του; Με ποιον τρόπο μπορούμε να αποδείξουμε τη σχέση αυτή; (σελ. 263)

      \item Με ποιον τρόπο μπορούμε να υπολογίσουμε τη χρονική στιγμή που ένα μέγεθος, το οποίο μεταβάλλεται με το χρόνο, αποκτά μια συγκεκριμένη τιμή; Ποιες τιμές θα λάβουμε υπόψη; (σελ. 267)

      \item Να περιγράψεις πώς μπορούμε να βρούμε τη χρονική διάρκεια μετάβασης ενός σώματος που εκτελεί απλή αρμονική ταλάντωση από μια θέση σε μια άλλη με τη βοήθεια του περιστρεφόμενου διανύσματος. (σελ. 270)

      \item Ποιου νόμου αποτελεί έκφραση ο τύπος $\Sigma F = -Dx$ στην απλή αρμονική ταλάντωση; (σελ. 273).

      \item Από ποιον τύπο πρέπει να ξεκινάς για να βρεις τη χρονική εξίσωση της συνισταμένης δύναμης που δέχεται ένα σώμα το οποίο εκτελεί απλή αρμονική ταλάντωση; (σελ. 273)

      \item Ποια είναι η μορφή της γραφικής παράστασης του τύπου $\Sigma F = -Dx$ και ποια η μορφή της γραφικής παράστασης για τον τύπο $a = -\omega^2 x$; (σελ. 273)

      \item Τι πληροφορίες μπορούμε να εξαγουμε από τη γραφική παράσταση της χρονικής εξίσωσης ενός μεγέθους για την εύρεση της αρχικής φάσης στην απλή αρμονική ταλάντωση; Ποιον άξονα πρέπει να παρατηρήσουμε ώστε να βγάλουμε ασφαλή συμπεράσματα για την περίοδο της ταλάντωσης; (σελ. 276)

      \item Πώς προκύπτουν οι χρονικές εξισώσεις της κινητικής ενέργειας του σώματος και της δυναμικής ενέργειας της ταλάντωσης σε μια απλή αρμονική ταλάντωση; (σελ. 279)

      \item Γιατί είναι πολύ χρήσιμος ο τύπος που προκύπτει εφαρμόζοντας την Αρχή Διατήρησης της Ενέργειας σε μια απλή αρμονική ταλάντωση; (μεθοδολογία + να μην ξεχάσεις, σελ. 279)

      \item Ποια είναι η συχνότητα μεγιστοποίησης και η συχνότητα μηδενισμού της κινητικής ενέργειας του σώματος και της δυναμικής ενέργειας της ταλάντωσης σε μια απλή αρμονική ταλάντωση; Μπορείς να αποδείξεις τον ισχυρισμό σου; (σελ. 284)

      \item Τι είναι αποδοτικό να χρησιμοποιούμε όταν μας δίνουν ή μας ζητούν σχέση ενεργειών; (σελ. 287)

      \item Με ποιους τρόπους μπορούμε να υπολογίσουμε το έργο της δύναμης επαναφοράς; Πώς συμπεριφέρεται ενεργειακά η δύναμη επαναφοράς; Εξαρτάται το έργο της δύναμης επαναφοράς από τη διαδρομή που ακολουθεί το σώμα μεταξύ της αρχικής και της τελικής θέσης; (σελ. 287)

      \item Η μορφή των εξισώσεων της κινητικής ενέργειας του σώματος και της δυναμικής ενέργειας της ταλάντωσης σε συνάρτηση με την απομάκρυνση ή την ταχύτητα σε μια απλή αρμονική ταλάντωση καθώς επίσης και οι αντίστοιχες γραφικές παραστάσεις εξαρτώνται από την αρχική φάση της ταλάντωσης; Ισχύει το ίδιο και για τις χρονικές εξισώσεις της κινητικής και της δυναμικής ενέργειας της ταλάντωσης; (σελ. 287)

      \item Με ποιον τύπο μπορούμε να υπολογίσουμε το ρυθμό μεταβολής της κινητικής ενέργειας ενός σώματος που εκτελεί απλή αρμονική ταλάντωση; Μπορείς να αποδείξεις τον τύπο αυτό; (σελ. 290 για την απόδειξη: 289)

      \item Με ποιον τύπο μπορούμε να υπολογίσουμε το ρυθμό μεταβολής της δυναμικής ενέργειας σε μια απλή αρμονική ταλάντωση; Μπορείς να αποδείξεις τον τύπο αυτό; (σελ. 290)
\end{enumerate}

\subsection{Σύστημα Μάζα - Ελατήριο, ΑΔΟ}

\begin{enumerate}
      \item Όταν φτιάχνεις σχήμα και πρέπει να μελετήσεις την κίνηση ενός σώματος μετά την κρούση, σε ποιο σημείο έχει σημασία να σχεδιάζεις το σώμα του οποίου την κίνηση θέλεις να μελετήσεις μετά την κρούση; (σελ. 415)

      \item Αλλάζει η θέση ισορροπίας της ταλάντωσης ενός σώματος ύστερα από μια πλαστική κρούση, αν το ελατήριο είναι οριζόντιο; Να αιτιολογήσεις την απάντησή σου. (σελ. 415)

      \item Αν έχεις ένα οριζόντιο ελατήριο και η κρούση συμβεί στη θέση φυσικού μήκους του ελατηρίου, τότε η ταχύτητα που αποκτά το σώμα, που είναι δεμένο με το ελατήριο και το οποίο θα εκτελέσει απλή αρμονική ταλάντωση, είναι η μέγιστη ταχύτητα της ταλάντωσης; Να αιτιολογήσεις την απάντησή σου. (σελ. 415)

      \item Σε ποια περίπτωση μετά από μια κρούση αλλάζει η θέση ισορροπίας της ταλάντωσης; (σελ. 419)

      \item Ποια απόσταση πρέπει να βρίσκεις όταν αλλάζει η θέση ισορροπίας; Πώς μπορείς να βρεις την απόσταση αυτή; (σελ. 420)

      \item Αν διαπιστώσεις ότι τη χρονική στιγμή που ξεκινά η απλή αρμονική ταλάντωση του σώματος που είναι δεμένο στο ελατήριο το σώμα βρίσκεται σε ενδιάμεση θέση της ταλάντωσής του, ποιος είναι ένας χρήσιμος τρόπος για να βρίσκεις το πλάτος της ταλάντωσης; (σελ. 424)

      \item Τι συμβαίνει με τη χρονική διάρκεια κίνησης δύο σωμάτων που ξεκινούν την κίνησή τους ταυτόχρονα από διαφορετικά σημεία και συναντιούνται; (σελ. 426)

      \item Αν έχεις πλαστική κρούση και κατακόρυφο ελατήριο σε κεκλιμένο επίπεδο, τι συμβαίνει με τη θέση ισορροπίας της ταλάντωσης; (σελ. 430)

      \item Πώς μπορείς να βρεις τη συνάρτηση της αλγεβρικής τιμής της δύναμης του ελατηρίου ως προς το χρόνο; (σελ. 432)

      \item Μπορείς να εφαρμόσεις την Α.Δ.Ο. σε μια έκρηξη; Να αιτιολογήσεις την απάντησή σου. (σελ. 441)
\end{enumerate}


\subsection{Σύστημα Μάζα - Ελατήριο, Χάσιμο Επαφής}

\begin{enumerate}
      \item Όταν φτιάχνεις σχήμα και πρέπει να μελετήσεις την κίνηση ενός σώματος μετά την κρούση, σε ποιο σημείο έχει σημασία να σχεδιάζεις το σώμα του οποίου την κίνηση θέλεις να μελετήσεις μετά την κρούση; (σελ. 415)

      \item Αλλάζει η θέση ισορροπίας της ταλάντωσης ενός σώματος ύστερα από μια πλαστική κρούση, αν το ελατήριο είναι οριζόντιο; Να αιτιολογήσεις την απάντησή σου. (σελ. 415)

      \item Αν έχεις ένα οριζόντιο ελατήριο και η κρούση συμβεί στη θέση φυσικού μήκους του ελατηρίου, τότε η ταχύτητα που αποκτά το σώμα, που είναι δεμένο με το ελατήριο και το οποίο θα εκτελέσει απλή αρμονική ταλάντωση, είναι η μέγιστη ταχύτητα της ταλάντωσης; Να αιτιολογήσεις την απάντησή σου. (σελ. 415)

      \item Σε ποια περίπτωση μετά από μια κρούση αλλάζει η θέση ισορροπίας της ταλάντωσης; (σελ. 419)

      \item Ποια απόσταση πρέπει να βρίσκεις όταν αλλάζει η θέση ισορροπίας; Πώς μπορείς να βρεις την απόσταση αυτή; (σελ. 420)

      \item Αν διαπιστώσεις ότι τη χρονική στιγμή που ξεκινά η απλή αρμονική ταλάντωση του σώματος που είναι δεμένο στο ελατήριο το σώμα βρίσκεται σε ενδιάμεση θέση της ταλάντωσής του, ποιος είναι ένας χρήσιμος τρόπος για να βρίσκεις το πλάτος της ταλάντωσης; (σελ. 424)

      \item Τι συμβαίνει με τη χρονική διάρκεια κίνησης δύο σωμάτων που ξεκινούν την κίνησή τους ταυτόχρονα από διαφορετικά σημεία και συναντιούνται; (σελ. 426)

      \item Αν έχεις πλαστική κρούση και κατακόρυφο ελατήριο σε κεκλιμένο επίπεδο, τι συμβαίνει με τη θέση ισορροπίας της ταλάντωσης; (σελ. 430)

      \item Πώς μπορείς να βρεις τη συνάρτηση της αλγεβρικής τιμής της δύναμης του ελατηρίου ως προς το χρόνο; (σελ. 432)

      \item Μπορείς να εφαρμόσεις την Α.Δ.Ο. σε μια έκρηξη; Να αιτιολογήσεις την απάντησή σου. (σελ. 441)
\end{enumerate}

\subsection{Φθίνουσες Μηχανικές Ταλαντώσεις}

\begin{enumerate}
      \item Ποιες είναι οι χρήσιμες ιδιότητες των λογαρίθμων που χρειάζεται να γνωρίζεις, ώστε να χειριστείς την εξίσωση του πλάτους στη φθίνουσα ταλάντωση; (σελ. 520)

      \item Τι σημαίνει χρόνος ημιζωής του πλάτους της φθίνουσας ταλάντωσης και ποια είναι η τιμή του; Ποια ενδιαφέρουσα ιδιότητα παρουσιάζει το πλάτος αναφορικά με το χρόνο ημιζωής; (σελ. 524)

      \item Ποια είναι η σχέση που ικανοποιούν τα πλάτα στο τέλος κάθε περιόδου της φθίνουσας ταλάντωσης; (σελ. 525)

      \item Με ποιον τύπο μπορούμε να υπολογίσουμε την ενέργεια ενός ταλαντωτή που εκτελεί φθίνουσα μηχανική ταλάντωση σε κάθε περίπτωση; Με ποιον τύπο μπορούμε να βρούμε προσεγγιστικά την ενέργεια του ταλαντωτή όταν η σταθερά απόσβεσης είναι μικρή; (σελ. 527)

      \item Με ποιον τύπο μπορούμε να υπολογίσουμε την απώλεια ενέργειας σε μια φθίνουσα ταλάντωση; (σελ. 527)

      \item Ποιος είναι ο τύπος της χρονικής εξίσωσης της ενέργειας της φθίνουσας ταλάντωσης με ικανοποιητική προσέγγιση όταν η σταθερά απόσβεσης είναι μικρή; (σελ. 529)

      \item Ποια σχέση συνδέει την απώλεια ενέργειας στη φθίνουσα ταλάντωση με το έργο της δύναμης αντίστασης στην κίνηση; (σελ. 531)

      \item Ποια σχέση προκύπτει από την εφαρμογή του 2ου νόμου του Νεύτωνα σε μια φθίνουσα ταλάντωση κατά την οποία το σώμα δέχεται τη δύναμη επαναφοράς και τη δύναμη αντίστασης στην κίνηση; (σελ. 533)

      \item Πώς μπορούμε να υπολογίσουμε το ρυθμό μεταβολής της ενέργειας μιας φθίνουσας ταλάντωσης; (σελ. 533)
\end{enumerate}

\subsection{Εξαναγκασμένες Ταλαντώσεις}

\begin{enumerate}
      \item Διαθέτεις ένα σύστημα μάζα - ιδανικό ελατήριο. Με ποιον τύπο μπορούμε να υπολογίσουμε την ιδιοσυχνότητα του συστήματος αυτού; (σελ. 560)

      \item Ποιες είναι οι τρεις συχνότητες που συναντάμε στη μελέτη των εξαναγκασμένων ταλαντώσεων; Ποιες είναι οι δύο συχνότητες οι οποίες είναι πάντοτε ίσες μεταξύ τους; (σελ. 560)

      \item Ποιες είναι οι τρεις δυνάμεις που ασκούνται στη διεύθυνση της κίνησης ενός σώματος που εκτελεί εξαναγκασμένη μηχανική ταλάντωση; Ποια είναι η σχέση στην οποία καταλήγουμε αν εφαρμόσουμε το 2ο νόμο του Νεύτωνα για μια τυχαία θέση της ταλάντωσης; (σελ. 560)

      \item Κατά τη διάρκεια της εξαναγκασμένης ταλάντωσης η μέγιστη κινητική ενέργεια είναι ίση με τη μέγιστη δυναμική ενέργεια της ταλάντωσης σε οποιαδήποτε περίπτωση; (σελ. 560)

      \item Ποια είναι η μορφή των χρονικών εξισώσεων της απομάκρυνσης, της ταχύτητας και της επιτάχυνσης στην εξαναγκασμένη ταλάντωση; (σελ. 560)

      \item Πότε ο ρυθμός με τον οποίο ο διεγέρτης προσφέρει ενέργεια στο σύστημα ισούται με το ρυθμό με τον οποίο η δύναμη αντίστασης στην κίνηση αφαιρεί ενέργεια από το σύστημα; Τι συμπέρασμα μπορείς να εξαγεις για τη μηχανική ενέργεια; (σελ. 563-564)

      \item Ποιες είναι οι σχέσεις που συνδυάζουν την απομάκρυνση $x$ με την επιτάχυνση $a$ καθώς και την απομάκρυνση $x$ με την ταχύτητα $u$ σε μια εξαναγκασμένη ταλάντωση; Πώς αποδεικνύονται οι σχέσεις αυτές; (σελ. 564)
\end{enumerate}

\section{Κύματα}

\begin{enumerate}
      \item Τι ονομάζουμε κύμα; Ένα κύμα μεταφέρει ενέργεια, ορμή και ύλη; (σελ. 13)

      \item Ποιες είναι οι κατηγορίες των κυμάτων ανάλογα με το είδος της ενέργειας που μεταφέρουν; (σελ. 13)

      \item Ποια κύματα ονομάζονται διαμήκη και ποια εγκάρσια; Σε ποια ελαστικά σώματα διαδίδονται τα εγκάρσια και σε ποια τα διαμήκη κύματα; (σελ. 13)

      \item Ποια κύματα δημιουργούν στο ελαστικό μέσο πυκνώματα και αραιώματα και ποια όρη και κοιλάδες; (σελ. 13)

      \item Σε ποιες κατηγορίες διακρίνονται τα κύματα ανάλογα με τις διαστάσεις του χώρου στον οποίο διαδίδονται; (σελ. 14)

      \item Από τι εξαρτάται η ταχύτητα διάδοσης ενός κύματος που διαδίδεται σε ελαστικό μέσο και από ποιον τύπο μπορεί να υπολογιστεί; (σελ. 14)

      \item Πότε ένα κύμα ονομάζεται αρμονικό; (σελ. 14)

      \item Τι ονομάζεται περίοδος, συχνότητα και μήκος κύματος ενός αρμονικού κύματος; (σελ. 14-15)

      \item Να γράψετε και να αποδείξετε τη θεμελιώδη εξίσωση της κυματικής. (σελ. 15)

      \item Ένα κύμα πλάτους $A$ διαδίδεται κατά μήκος γραμμικού ελαστικού μέσου που ταυτίζεται με τον άξονα $x'Ox$ προς τη θετική κατεύθυνση του άξονα και τη χρονική στιγμή $t = 0$ το κύμα φτάνει στην αρχή $O(x = 0)$ του άξονα θέτοντας σε απλή αρμονική ταλάντωση το υλικό σημείο που βρίσκεται στην αρχή $O$ με ταχύτητα θετικής φοράς. Να γράψετε την εξίσωση του αρμονικού κύματος και να την αποδείξετε. (σελ. 15-16)

      \item Ένα κύμα πλάτους $A$ διαδίδεται κατά μήκος γραμμικού ελαστικού μέσου που ταυτίζεται με τον άξονα $x'Ox$ προς την αρνητική κατεύθυνση του άξονα και τη χρονική στιγμή $t = 0$ το κύμα φτάνει στην αρχή $O(x = 0)$ του άξονα θέτοντας σε απλή αρμονική ταλάντωση το υλικό σημείο που βρίσκεται στην αρχή $O$ με ταχύτητα θετικής φοράς. Να γράψετε την εξίσωση του αρμονικού κύματος και να την αποδείξετε. (σελ. 17)

      \item Τι ονομάζεται φάση ενός αρμονικού κύματος; Τι τιμές μπορεί να πάρει η φάση ενός κύματος; (σελ. 17)

      \item Ένα αρμονικό κύμα διαδίδεται κατά μήκος γραμμικού ελαστικού μέσου που ταυτίζεται με τον άξονα $x'Ox$ και η εξίσωσή του είναι της μορφής $y = A\eta\mu 2\pi (\frac{t}{T} \pm \frac{x}{\lambda})$. Ένα υλικό σημείο $M(x_M)$ του μέσου ξεκινά να ταλαντώνεται τη χρονική στιγμή $t_M$. Να γράψετε τις χρονικές εξισώσεις της απομάκρυνσης $y_M = f(t)$ από τη θέση ισορροπίας, της ταχύτητας ταλάντωσης $u_M = f(t)$ και της επιτάχυνσης ταλάντωσης $a_M = f(t)$ του υλικού σημείου $M$. (σελ. 17-18)

      \item Ποια γραφική παράσταση ονομάζεται στιγμιότυπο του κύματος; (σελ. 18-19)

      \item Να αποδείξετε ότι η φάση ενός αρμονικού κύματος μια συγκεκριμένη χρονική στιγμή $t_1$ για ορισμένη θέση $x_1$, είναι αδύνατον να είναι αρνητική. (σελ. 19)

      \item Να σχεδιάσετε ποιοτικά τη γραφική παράσταση της φάσης ενός αρμονικού κύματος σε συνάρτηση με τις θέσεις $x$ των υλικών σημείων πάνω στον άξονα $x'Ox$ για ένα κύμα που διαδίδεται: i) προς τα δεξιά, ii) προς τ' αριστερά. Με τι ισούται η κλίση κάθε γραφικής παράστασης που σχεδιάσατε; (σελ. 19-20)

      \item Να σχεδιάσετε ποιοτικά τη γραφική παράσταση της φάσης της ταλάντωσης ενός υλικού σημείου που βρίσκεται στη θέση $x_1$ του ελαστικού μέσου, στο οποίο έφτασε το αρμονικό κύμα μετά τη χρονική στιγμή $t = 0$, σε συνάρτηση με το χρόνο $t$. Με τι ισούται η κλίση της γραφικής παράστασης που σχεδιάσατε; (σελ. 20-21)

      \item Να γράψετε και να αποδείξετε τον τύπο υπολογισμού της διαφοράς φάσης μεταξύ των ταλαντώσεων δύο υλικών σημείων του ελαστικού μέσου στο οποίο διαδίδεται ένα αρμονικό κύμα την ίδια χρονική στιγμή. (σελ. 21)

      \item Πότε λέμε ότι δύο υλικά σημεία του ελαστικού μέσου στο οποίο διαδίδεται ένα αρμονικό κύμα βρίσκονται σε συμφωνία και πότε σε αντίθεση φάσης; (σελ. 22)

      \item Με ποιον τύπο μπορούμε να υπολογίσουμε τη μεταβολή της φάσης της ταλάντωσης ενός υλικού σημείου μεταξύ δύο χρονικών στιγμών $t_1$ και $t_2$ ($t_2 > t_1$) και πώς αποδεικνύεται ο τύπος αυτός; (σελ. 22-23)

      \item Τι σημαίνει ότι ένα αρμονικό κύμα έχει αρχική φάση; (σελ. 23)
\end{enumerate}

\begin{enumerate}
      \item Πώς διατυπώνεται η αρχή της επαλληλίας (ή υπέρθεσης) κυμάτων; (σελ. 23-24)

      \item Ποια συμπεράσματα προκύπτουν από την επαλληλία δύο κυμάτων που διαδίδονται στο ίδιο ελαστικό μέσο; (σελ. 24)

      \item Πότε παραβιάζεται η αρχή της επαλληλίας; (σελ. 24)

      \item Ποιες πηγές κυμάτων ονομάζονται σύγχρονες και ποιες σύμφωνες; (σελ. 24)

      \item Ποιο φαινόμενο ονομάζεται συμβολή; (σελ. 24)

      \item Πότε λέμε ότι σ' ένα υλικό σημείο του μέσου διάδοσης συμβαίνει ενισχυτική συμβολή και πότε ακυρωτική (ή αποσβεστική) συμβολή; (σελ. 25)

      \item Σε δύο σημεία της επιφάνειας ενός ελαστικού μέσου υπάρχουν δύο σύγχρονες πηγές αρμονικών κυμάτων που ταλαντώνονται με εξίσωση $y_1 = y_2 = A\eta\mu\frac{2\pi}{T}t$. Να αποδείξετε τη χρονική εξίσωση ταλάντωσης ενός υλικού σημείου $\Sigma$ του μέσου διάδοσης από τη στιγμή που συμβάλλουν και τα δύο κύματα στο σημείο αυτό. (σελ. 25-26) - \textbf{εκτός ύλης}

      \item Σε υλικό σημείο $\Sigma$ της επιφάνειας του ελαστικού μέσου, που απέχει απόσταση $r_1$ από την πηγή $\Pi_1$ και απόσταση $r_2$ από την πηγή $\Pi_2$, συμβαίνει ενισχυτική συμβολή. Ποια σχέση ικανοποιούν οι αποστάσεις $r_1$ και $r_2$; (σελ. 26) - \textbf{η απόδειξη εκτός ύλης}

      \item Σε υλικό σημείο $Z$ της επιφάνειας του ελαστικού μέσου, που απέχει απόσταση $r_1$ από την πηγή $\Pi_1$ και απόσταση $r_2$ από την πηγή $\Pi_2$, συμβαίνει ακυρωτική (ή αποσβεστική) συμβολή. Ποια σχέση ικανοποιούν οι αποστάσεις $r_1$ και $r_2$; (σελ. 26-27) - \textbf{η απόδειξη εκτός ύλης}

      \item Δύο σύγχρονες πηγές $\Pi_1$ και $\Pi_2$ δημιουργούν στην επιφάνεια ενός ελαστικού μέσου αρμονικά κύματα τα οποία συμβάλλουν. Να σχεδιάσετε το γεωμετρικό τόπο των σημείων του μέσου που εμφανίζουν ενισχυτική και ακυρωτική συμβολή (κροσσοί συμβολής). (σελ. 27)
\end{enumerate}

\begin{enumerate}
      \item Τι ονομάζεται στάσιμο κύμα και πώς δημιουργείται; (σελ. 27)

      \item Δύο αρμονικά κύματα διαδίδονται στο ίδιο γραμμικό ελαστικό μέσο, που ταυτίζεται με τον άξονα $x'Ox$, έχουν εξισώσεις $y_1 = A\eta\mu 2\pi (\frac{t}{T} - \frac{x}{\lambda})$ και $y_2 = A\eta\mu 2\pi (\frac{t}{T} + \frac{x}{\lambda})$ και λόγω της συμβολής τους δημιουργείται στάσιμο κύμα. Να γράψετε και να αποδείξετε την εξίσωση του στάσιμου κύματος. (σελ. 27-28)

      \item Σε γραμμικό ελαστικό μέσο, που ταυτίζεται με τον άξονα $x'Ox$, έχει δημιουργηθεί στάσιμο κύμα με κοιλία στην αρχή μέτρησης των αποστάσεων ($x = 0$). Να γράψετε και να αποδείξετε τη σχέση υπολογισμού των θέσεων του άξονα όπου τα υλικά σημεία που βρίσκονται στις θέσεις αυτές, ταλαντώνονται με μέγιστο πλάτος (κοιλίες). (σελ. 28-29)

      \item Σε γραμμικό ελαστικό μέσο, που ταυτίζεται με τον άξονα $x'Ox$, έχει δημιουργηθεί στάσιμο κύμα με κοιλία στην αρχή μέτρησης των αποστάσεων ($x = 0$). Να γράψετε και να αποδείξετε τη σχέση υπολογισμού των θέσεων του άξονα όπου τα υλικά σημεία που βρίσκονται στις θέσεις αυτές παραμένουν συνεχώς ακίνητα (δεσμοί). (σελ. 29)

      \item Σε γραμμικό ελαστικό μέσο έχει δημιουργηθεί στάσιμο κύμα εξαιτίας της ταυτόχρονης διάδοσης δύο αρμονικών κυμάτων που έχουν ίδιο πλάτος $A$, ίδια συχνότητα $f$ και ίδιο μήκος κύματος $\lambda$, ενώ διαδίδονται με αντίθετη φορά. Να γράψετε και να αποδείξετε με τι ισούται: i) η απόσταση μεταξύ δύο διαδοχικών κοιλιών, ii) η απόσταση μεταξύ δύο διαδοχικών δεσμών, iii) η απόσταση μεταξύ ενός δεσμού και της πλησιέστερης σε αυτόν κοιλίας. (σελ. 30)

      \item Ποια είναι η διαφορά φάσης μεταξύ των ταλαντώσεων δύο υλικών σημείων του ελαστικού μέσου στο οποίο δημιουργείται στάσιμο κύμα; (σελ. 30)
\end{enumerate}

\begin{enumerate}
      \item Τι είναι το ηλεκτρομαγνητικό κύμα; (σελ. 31)

      \item Πώς παράγεται το ηλεκτρομαγνητικό κύμα με τη βοήθεια ενός ταλαντευόμενου ηλεκτρικού διπόλου; (σελ. 31)

      \item Ποια συμπεράσματα προκύπτουν από τη μελέτη των ηλεκτρομαγνητικών κυμάτων; (σελ. 32)

      \item Μπορεί να δημιουργηθεί ηλεκτρομαγνητικό κύμα από ακίνητα ηλεκτρικά φορτία ή φορτία που κινούνται με σταθερή ταχύτητα; (σελ. 32)

      \item Ποιες είναι οι εξισώσεις που περιγράφουν ένα ηλεκτρομαγνητικό κύμα; (σελ. 32)

      \item Ποιοι είναι οι τρόποι δημιουργίας ηλεκτρομαγνητικών κυμάτων; (σελ. 33)

      \item Ποιες είναι οι κατηγορίες στις οποίες χωρίζουμε τα ηλεκτρομαγνητικά κύματα (φάσμα της ηλεκτρομαγνητικής ακτινοβολίας); (σελ. 33-34)

      \item Τι είναι το ορατό φως; (σελ. 34)

      \item Ποια μήκη κύματος του φωτός προκαλούν την αίσθηση συγκεκριμένων χρωμάτων; (σελ. 34)

      \item Ποια ακτινοβολία ονομάζεται μονοχρωματική; (σελ. 34)
\end{enumerate}

\subsection{Φάση Αρμονικού Κύματος}

\begin{enumerate}
      \item Ποιος είναι ο τύπος της φάσης ενός αρμονικού κύματος το οποίο διαδίδεται προς τη θετική κατεύθυνση και ποιος όταν διαδίδεται προς την αρνητική κατεύθυνση; (σελ. 320)

      \item Ποιες είναι οι δύο ανεξάρτητες μεταβλητές από τις οποίες εξαρτάται η φάση ενός αρμονικού κύματος; (σελ. 320)

      \item Γιατί η φάση ενός αρμονικού κύματος δεν μπορεί να πάρει αρνητικές τιμές; (σελ. 320)

      \item Ποια είναι η φυσική σημασία της φάσης ενός αρμονικού κύματος; (σελ. 321)

      \item Τι ονομάζουμε μέτωπο κύματος και ποια είναι η φάση των σημείων που βρίσκονται πάνω σε ένα μέτωπο κύματος; (σελ. 321)

      \item Πώς μεταβάλλεται η φάση της ταλάντωσης ενός υλικού σημείου $x_1$ του ελαστικού μέσου σε συνάρτηση με το χρόνο; Να γράψετε τη σχετική εξίσωση. (σελ. 322)

      \item Να σχεδιάσετε τη γραφική παράσταση της φάσης της ταλάντωσης ενός υλικού σημείου $x_1$ σε συνάρτηση με το χρόνο, για την περίπτωση που το κύμα διαδίδεται προς τη θετική κατεύθυνση. (σελ. 323)

      \item Τι παριστάνει η κλίση της γραφικής παράστασης της φάσης σε συνάρτηση με το χρόνο για ένα υλικό σημείο $x_1$; (σελ. 323)

      \item Ποια είναι η τιμή της φάσης ενός υλικού σημείου $x_1$ πριν φτάσει το κύμα σε αυτό; (σελ. 324)

      \item Πώς μεταβάλλεται η φάση των υλικών σημείων του ελαστικού μέσου μια συγκεκριμένη χρονική στιγμή $t_1$ σε συνάρτηση με τη θέση τους; Να γράψετε τη σχετική εξίσωση. (σελ. 324)

      \item Να σχεδιάσετε τη γραφική παράσταση της φάσης των υλικών σημείων του ελαστικού μέσου σε συνάρτηση με τη θέση τους, μια συγκεκριμένη χρονική στιγμή $t_1$, για την περίπτωση που το κύμα διαδίδεται προς τη θετική κατεύθυνση. (σελ. 325)

      \item Τι παριστάνει η κλίση της γραφικής παράστασης της φάσης σε συνάρτηση με τη θέση των υλικών σημείων μια συγκεκριμένη χρονική στιγμή $t_1$, για ένα κύμα που διαδίδεται προς τη θετική κατεύθυνση; (σελ. 325)

      \item Να σχεδιάσετε τη γραφική παράσταση της φάσης των υλικών σημείων του ελαστικού μέσου σε συνάρτηση με τη θέση τους, μια συγκεκριμένη χρονική στιγμή $t_1$, για την περίπτωση που το κύμα διαδίδεται προς την αρνητική κατεύθυνση. (σελ. 326)

      \item Τι παριστάνει η κλίση της γραφικής παράστασης της φάσης σε συνάρτηση με τη θέση των υλικών σημείων μια συγκεκριμένη χρονική στιγμή $t_1$, για ένα κύμα που διαδίδεται προς την αρνητική κατεύθυνση; (σελ. 326)

      \item Μέχρι ποια θέση σχεδιάζουμε τη γραφική παράσταση της φάσης σε συνάρτηση με τη θέση των υλικών σημείων μια συγκεκριμένη χρονική στιγμή $t_1$; (σελ. 326)

      \item Με ποιον τύπο υπολογίζεται η διαφορά φάσης μεταξύ των ταλαντώσεων δύο υλικών σημείων $x_1$ και $x_2$ του ελαστικού μέσου την ίδια χρονική στιγμή; (σελ. 331)

      \item Η διαφορά φάσης μεταξύ των ταλαντώσεων δύο υλικών σημείων $x_1$ και $x_2$ του ελαστικού μέσου εξαρτάται από τη χρονική στιγμή που τη μελετάμε; (σελ. 331)

      \item Πότε δύο υλικά σημεία του ελαστικού μέσου βρίσκονται σε συμφωνία φάσης και ποια είναι η ελάχιστη απόσταση μεταξύ δύο τέτοιων σημείων; (σελ. 331)

      \item Πότε δύο υλικά σημεία του ελαστικού μέσου βρίσκονται σε αντίθεση φάσης και ποια είναι η ελάχιστη απόσταση μεταξύ δύο τέτοιων σημείων; (σελ. 332)

      \item Με ποιον τύπο υπολογίζεται η μεταβολή της φάσης της ταλάντωσης ενός υλικού σημείου $x_1$ μεταξύ δύο χρονικών στιγμών $t_1$ και $t_2$; (σελ. 338)

      \item Τι σημαίνει ότι ένα αρμονικό κύμα έχει αρχική φάση; (σελ. 343)

      \item Ποια είναι η διαφορά φάσης μεταξύ της ταχύτητας ταλάντωσης και της απομάκρυνσης ενός υλικού σημείου του ελαστικού μέσου; (σελ. 343)

      \item Πώς συνδέεται η έννοια της φάσης με το στιγμιότυπο ενός κύματος; (σελ. 343)
\end{enumerate}

\subsection{Συμβολή και Στάσιμα Κύματα}

\begin{enumerate}
      \item Τι ονομάζουμε εξίσωση ενός κύματος; (σελ. 301)

      \item Τι περιγράφει η χρονική εξίσωση κίνησης $y = f(t)$ ενός υλικού σημείου του ελαστικού μέσου στο οποίο διαδίδεται ένα κύμα; (σελ. 302)

      \item Ποια είναι η μορφή της χρονικής εξίσωσης της ταχύτητας ταλάντωσης και ποια της επιτάχυνσης ταλάντωσης ενός υλικού σημείου του ελαστικού μέσου στο οποίο διαδίδεται ένα αρμονικό κύμα; (σελ. 302-303)

      \item Ποια είναι η διαφορά μεταξύ της ταχύτητας διάδοσης $v$ ενός κύματος και της ταχύτητας ταλάντωσης $u$ των υλικών σημείων του ελαστικού μέσου; (σελ. 303)

      \item Τι περιγράφει η εξίσωση $y = f(x)$ η οποία μας δίνει τις απομακρύνσεις όλων των υλικών σημείων του ελαστικού μέσου από τη θέση ισορροπίας τους μια συγκεκριμένη χρονική στιγμή $t_1$; (σελ. 304)

      \item Πώς μεταβάλλεται η ταχύτητα ταλάντωσης των υλικών σημείων ενός ελαστικού μέσου σε συνάρτηση με τη θέση τους $x$, μια συγκεκριμένη χρονική στιγμή $t_1$, κατά τη διάδοση ενός αρμονικού κύματος; Να γράψετε τη σχετική εξίσωση. (σελ. 305)

      \item Από ποια σχέση προκύπτει η εξίσωση ενός αρμονικού κύματος $y = A\eta\mu 2\pi (\frac{t}{T} - \frac{x}{\lambda})$ το οποίο διαδίδεται κατά τη θετική κατεύθυνση; (σελ. 306)

      \item Από ποια σχέση προκύπτει η εξίσωση ενός αρμονικού κύματος $y = A\eta\mu 2\pi (\frac{t}{T} + \frac{x}{\lambda})$ το οποίο διαδίδεται κατά την αρνητική κατεύθυνση; (σελ. 308)

      \item Ποιες είναι οι εναλλακτικές μορφές με τις οποίες μπορεί να γραφεί η εξίσωση ενός αρμονικού κύματος; (σελ. 312)

      \item Πώς μπορούμε να διαπιστώσουμε αν ένα κύμα, του οποίου μας δίνουν την εξίσωση, διαδίδεται προς τη θετική ή προς την αρνητική κατεύθυνση του άξονα; (σελ. 312)

      \item Ποια είναι η χρονική εξίσωση της ταχύτητας ταλάντωσης $u = f(t,x)$ και ποια της επιτάχυνσης ταλάντωσης $a = f(t,x)$ ενός τυχαίου υλικού σημείου $x$ του ελαστικού μέσου, τη χρονική στιγμή $t$, για ένα αρμονικό κύμα που διαδίδεται προς τη θετική κατεύθυνση; (σελ. 314)

      \item Ποια είναι η χρονική εξίσωση της ταχύτητας ταλάντωσης $u = f(t,x)$ και ποια της επιτάχυνσης ταλάντωσης $a = f(t,x)$ ενός τυχαίου υλικού σημείου $x$ του ελαστικού μέσου, τη χρονική στιγμή $t$, για ένα αρμονικό κύμα που διαδίδεται προς την αρνητική κατεύθυνση; (σελ. 314)

      \item Με ποιον τρόπο μπορούμε να προσδιορίσουμε την κατεύθυνση κίνησης (πάνω ή κάτω) ενός υλικού σημείου $x_1$ του ελαστικού μέσου μια συγκεκριμένη χρονική στιγμή $t_1$, αν γνωρίζουμε το στιγμιότυπο του κύματος; (σελ. 315)

      \item Με ποιον τρόπο μπορούμε να προσδιορίσουμε την κατεύθυνση διάδοσης ενός κύματος (δεξιά ή αριστερά), αν γνωρίζουμε το στιγμιότυπο του κύματος και την κατεύθυνση κίνησης ενός υλικού σημείου του ελαστικού μέσου; (σελ. 316)
\end{enumerate}

\section{Κβαντομηχανική}

\begin{enumerate}
      \item Τι προβλέπει η κλασική θεωρία για την ενέργεια που μεταφέρει η ηλεκτρομαγνητική ακτινοβολία και ποια φαινόμενα αποτυγχάνει να ερμηνεύσει; (σελ. 225)

      \item Ποια ριζοσπαστική υπόθεση έκανε το 1900 ο Max Planck, η οποία ήταν το θεμέλιο της κβαντικής θεωρίας; (σελ. 225)

      \item Τι φαινόμενα ερμηνεύει η κβαντική θεωρία και πότε τα αποτελέσματά της ταυτίζονται με αυτά της κλασικής θεωρίας; (σελ. 226)

      \item Τι είναι η θερμική ακτινοβολία των σωμάτων και τι εκφράζει το μέγεθος που ονομάζεται ένταση της ακτινοβολίας; (σελ. 226)

      \item Τι ονομάζεται μέλαν (μαύρο) σώμα στη φυσική; (σελ. 226-227)

      \item Τι ονομάζεται ακτινοβολία μέλανος σώματος; (σελ. 227)

      \item Τι αναφέρει ο νόμος μετατόπισης του Wien; (σελ. 228)

      \item Ποια είναι η κλασική ερμηνεία για την ακτινοβολία του μέλανος σώματος και γιατί η ερμηνεία αυτή αποτυγχάνει; (σελ. 228)

      \item Ποιες είναι οι δύο υποθέσεις του Planck με τις οποίες ερμηνεύεται η ακτινοβολία του μέλανος σώματος; (σελ. 229)

      \item Μπορεί η κβάντωση της ενέργειας να παρατηρηθεί σε έναν ταλαντωτή στο μακρόκοσμο; (σελ. 230)

      \item Τι ονομάζεται φωτοηλεκτρικό φαινόμενο; (σελ. 230)

      \item Με ποια διάταξη μπορούμε να μελετήσουμε το φωτοηλεκτρικό φαινόμενο; (σελ. 230-231)

      \item Τι διαπιστώνεται πειραματικά με τη βοήθεια της διάταξης της σελίδας 231; (σελ. 231)

      \item Ποια είναι η γραφική παράσταση της έντασης $i$ του ρεύματος σε συνάρτηση με τη διαφορά δυναμικού $\Delta V$ μεταξύ της ανόδου και της καθόδου για δεδομένη συχνότητα και ένταση της ακτινοβολίας; (σελ. 232-233)

      \item Τι παρατηρούμε από τη γραφική παράσταση της έντασης $i$ του ρεύματος σε συνάρτηση με τη διαφορά δυναμικού $\Delta V$ μεταξύ της ανόδου και της καθόδου για δεδομένη συχνότητα και ένταση της ακτινοβολίας και τι ονομάζεται τάση αποκοπής $V_0$; (σελ. 232)

      \item Ποια είναι η γραφική παράσταση της έντασης $i$ σε συνάρτηση με τη διαφορά δυναμικού $\Delta V$ μεταξύ της ανόδου και της καθόδου για δεδομένη συχνότητα αλλά διαφορετική ένταση της ακτινοβολίας; (σελ. 233)

      \item Ποια είναι τα σημεία αποτυχίας της κλασικής θεωρίας στην ερμηνεία του φωτοηλεκτρικού φαινομένου; (σελ. 233-234)

      \item Πώς ορίζεται το ένα ηλεκτρονιοβόλτ; (σελ. 234)

      \item Ποια υπόθεση έκανε ο Einstein το 1905 για να ερμηνεύσει τα πειραματικά δεδομένα του φωτοηλεκτρικού φαινομένου; (σελ. 235)

      \item Ποια είναι η σχέση που συνδέει την ορμή ενός φωτονίου μιας ακτινοβολίας με το μήκος κύματος της ακτινοβολίας; (σελ. 235)

      \item Τι ονομάζεται έργο εξαγωγής ηλεκτρονίων από την επιφάνεια ενός μετάλλου; (σελ. 235-236)

      \item Ποια είναι η φωτοηλεκτρική εξίσωση του Einstein και πώς η εξίσωση αυτή ερμηνεύει την ύπαρξη της συχνότητας κατωφλίου; (σελ. 236-237)

      \item Ποια είναι η γραφική παράσταση της κινητικής ενέργειας των φωτοηλεκτρονίων σε συνάρτηση με τη συχνότητα της προσπίπτουσας ακτινοβολίας για τιμές ίσες ή μεγαλύτερες από τη συχνότητα κατωφλίου; (σελ. 237-238)

      \item Ποια είναι η σχέση της τάσης αποκοπής με τη συχνότητα της προσπίπτουσας ακτινοβολίας και ποια η αντίστοιχη γραφική παράσταση; (σελ. 238-239)

      \item Με τι ισούται η ένταση της ακτινοβολίας που προσπίπτει στην κάθοδο, χρησιμοποιώντας την υπόθεση του Einstein; (σελ. 239)

      \item Για ποιο λόγο η αύξηση της έντασης της ακτινοβολίας για δεδομένη συχνότητα προκαλεί μεγαλύτερη ένταση ρεύματος στο κύκλωμα σύμφωνα με την κβαντική θεώρηση; (σελ. 239-240)

      \item Ποια ακτινοβολία ονομάζεται ακτίνες Χ και ποιος είναι ο μηχανισμός παραγωγής των ακτίνων Χ; (σελ. 240)

      \item Τι αναφέρει το φαινόμενο Compton; (σελ. 241)

      \item Γιατί η κλασική θεωρία αποτυγχάνει στην εξήγηση του φαινομένου Compton; (σελ. 241)

      \item Πώς εξηγείται το φαινόμενο Compton; (σελ. 241-242)

      \item Πότε η κινητική ενέργεια που αποκτά ακίνητο ηλεκτρόνιο μετά την ελαστική κρούση με φωτόνιο είναι μηδενική και πότε μέγιστη; (σελ. 242-243)

      \item Ποια είναι η υπόθεση του Louis de Broglie για την κυματική φύση της ύλης; (σελ. 243)

      \item Γιατί δεν μπορούμε να παρατηρήσουμε τις κυματικές ιδιότητες ενός αντικειμένου που κινείται στο μακρόκοσμο; (σελ. 244)

      \item Για ποιο λόγο η κβαντική θεώρηση ότι ένα σωματίδιο έχει και κυματικές ιδιότητες εισάγει αβεβαιότητα στη μέτρηση της θέσης του και της ορμής του; (σελ. 244)

      \item Πώς μπορούμε να περιγράψουμε ένα σωματίδιο που έχει και κυματικές ιδιότητες; (σελ. 244-245)

      \item Η ύπαρξη αβεβαιότητας στη μέτρηση της θέσης και της ορμής ενός σωματιδίου εξαρτάται από την πειραματική συσκευή που χρησιμοποιούμε; (σελ. 245-246)

      \item Ποια είναι η διατύπωση της αρχής της αβεβαιότητας του Heisenberg; (σελ. 246)

      \item Ποια είναι η αρχή της αβεβαιότητας χρόνου - ενέργειας; (σελ. 246)

      \item Πώς μπορεί να ερμηνευτεί το εύρος των φασματικών γραμμών με τη βοήθεια της αρχής της αβεβαιότητας χρόνου - ενέργειας; (σελ. 247-248)

      \item Η αρχή της αβεβαιότητας μπορεί να εφαρμοστεί για αντικείμενα του μακρόκοσμου; (σελ. 248)

      \item Ποια ερμηνεία έδωσε ο Max Born στην κυματοσυνάρτηση που περιγράφει το κύμα το οποίο είναι συνυφασμένο με ένα σωματίδιο σύμφωνα με την κβαντική θεωρία; (σελ. 248-249)

      \item Ποια είναι η συνθήκη κανονικοποίησης που πρέπει να ικανοποιεί η κυματοσυνάρτηση $\Psi$; (σελ. 249)
\end{enumerate}

\subsection{ΑΚΤΙΝΟΒΟΛΙΑ ΜΕΛΑΝΟΣ ΣΩΜΑΤΟΣ}
\begin{enumerate}
      \setcounter{enumi}{3}
      \item Τι είναι η θερμική ακτινοβολία των σωμάτων και τι εκφράζει το μέγεθος που ονομάζεται ένταση της ακτινοβολίας; (σελ. 226)

      \item Τι ονομάζεται μέλαν (μαύρο) σώμα στη φυσική; (σελ. 226-227)

      \item Τι ονομάζεται ακτινοβολία μέλανος σώματος; (σελ. 227)

      \item Τι αναφέρει ο νόμος μετατόπισης του Wien; (σελ. 228)

      \item Ποια είναι η κλασική ερμηνεία για την ακτινοβολία του μέλανος σώματος και γιατί η ερμηνεία αυτή αποτυγχάνει; (σελ. 228)

      \item Ποιες είναι οι δύο υποθέσεις του Planck με τις οποίες ερμηνεύεται η ακτινοβολία του μέλανος σώματος; (σελ. 229)

      \item Μπορεί η κβάντωση της ενέργειας να παρατηρηθεί σε έναν ταλαντωτή στο μακρόκοσμο; (σελ. 230)
\end{enumerate}

\subsection{ΦΩΤΟΗΛΕΚΤΡΙΚΟ ΦΑΙΝΟΜΕΝΟ}
\begin{enumerate}
      \setcounter{enumi}{10}
      \item Τι ονομάζεται φωτοηλεκτρικό φαινόμενο; (σελ. 230)

      \item Με ποια διάταξη μπορούμε να μελετήσουμε το φωτοηλεκτρικό φαινόμενο; (σελ. 230-231)

      \item Τι διαπιστώνεται πειραματικά με τη βοήθεια της διάταξης της σελίδας 231; (σελ. 231)

      \item Ποια είναι η γραφική παράσταση της έντασης $i$ του ρεύματος σε συνάρτηση με τη διαφορά δυναμικού $\Delta V$ μεταξύ της ανόδου και της καθόδου για δεδομένη συχνότητα και ένταση της ακτινοβολίας; (σελ. 232-233)

      \item Τι παρατηρούμε από τη γραφική παράσταση της έντασης $i$ του ρεύματος σε συνάρτηση με τη διαφορά δυναμικού $\Delta V$ μεταξύ της ανόδου και της καθόδου για δεδομένη συχνότητα και ένταση της ακτινοβολίας και τι ονομάζεται τάση αποκοπής $V_0$; (σελ. 232)

      \item Ποια είναι η γραφική παράσταση της έντασης $i$ σε συνάρτηση με τη διαφορά δυναμικού $\Delta V$ μεταξύ της ανόδου και της καθόδου για δεδομένη συχνότητα αλλά διαφορετική ένταση της ακτινοβολίας; (σελ. 233)

      \item Ποια είναι τα σημεία αποτυχίας της κλασικής θεωρίας στην ερμηνεία του φωτοηλεκτρικού φαινομένου; (σελ. 233-234)

      \item Πώς ορίζεται το ένα ηλεκτρονιοβόλτ; (σελ. 234)

      \item Ποια υπόθεση έκανε ο Einstein το 1905 για να ερμηνεύσει τα πειραματικά δεδομένα του φωτοηλεκτρικού φαινομένου; (σελ. 235)

      \item Ποια είναι η σχέση που συνδέει την ορμή ενός φωτονίου μιας ακτινοβολίας με το μήκος κύματος της ακτινοβολίας; (σελ. 235)

      \item Τι ονομάζεται έργο εξαγωγής ηλεκτρονίων από την επιφάνεια ενός μετάλλου; (σελ. 235-236)

      \item Ποια είναι η φωτοηλεκτρική εξίσωση του Einstein και πώς η εξίσωση αυτή ερμηνεύει την ύπαρξη της συχνότητας κατωφλίου; (σελ. 236-237)

      \item Ποια είναι η γραφική παράσταση της κινητικής ενέργειας των φωτοηλεκτρονίων σε συνάρτηση με τη συχνότητα της προσπίπτουσας ακτινοβολίας για τιμές ίσες ή μεγαλύτερες από τη συχνότητα κατωφλίου; (σελ. 237-238)

      \item Ποια είναι η σχέση της τάσης αποκοπής με τη συχνότητα της προσπίπτουσας ακτινοβολίας και ποια η αντίστοιχη γραφική παράσταση; (σελ. 238-239)

      \item Με τι ισούται η ένταση της ακτινοβολίας που προσπίπτει στην κάθοδο, χρησιμοποιώντας την υπόθεση του Einstein; (σελ. 239)

      \item Για ποιο λόγο η αύξηση της έντασης της ακτινοβολίας για δεδομένη συχνότητα προκαλεί μεγαλύτερη ένταση ρεύματος στο κύκλωμα σύμφωνα με την κβαντική θεώρηση; (σελ. 239-240)
\end{enumerate}

\subsection{ΦΑΙΝΟΜΕΝΟ COMPTON}
\begin{enumerate}
      \setcounter{enumi}{26}
      \item Ποια ακτινοβολία ονομάζεται ακτίνες Χ και ποιος είναι ο μηχανισμός παραγωγής των ακτίνων Χ; (σελ. 240)

      \item Τι αναφέρει το φαινόμενο Compton; (σελ. 241)

      \item Γιατί η κλασική θεωρία αποτυγχάνει στην εξήγηση του φαινομένου Compton; (σελ. 241)

      \item Πώς εξηγείται το φαινόμενο Compton; (σελ. 241-242)

      \item Ποια είναι η σχέση που δίνει τη μεταβολή του μήκους κύματος $\Delta\lambda$ του σκεδαζόμενου φωτονίου σε συνάρτηση με τη γωνία σκέδασης $\theta$ (τύπος του Compton) και τι ονομάζεται μήκος κύματος Compton; (σελ. 242)

      \item Πότε η μεταβολή του μήκους κύματος $\Delta\lambda$ του φωτονίου κατά το φαινόμενο Compton είναι μηδενική και πότε παίρνει τη μέγιστη τιμή της; (σελ. 242)

      \item Πότε η κινητική ενέργεια που αποκτά το αρχικά ακίνητο ηλεκτρόνιο μετά την ελαστική κρούση με το φωτόνιο είναι μηδενική και πότε μέγιστη; (σελ. 242-243)

      \item Γιατί το φαινόμενο Compton δεν γίνεται αντιληπτό όταν οι ακτίνες Χ προσπίπτουν σε ηλεκτρόνια που είναι ισχυρά δεσμευμένα στα άτομα του στόχου; (σελ. 243)
\end{enumerate}

\section{Στερεό}

\subsection{Θεωρία μηχανικής στερεού σώματος}

\begin{enumerate}
      \item Τι ονομάζεται υλικό σημείο και τι μηχανικό στερεό σώμα; (σελ. 11)
      \item Πότε ένα στερεό σώμα εκτελεί μεταφορική κίνηση και πότε περιστροφική κίνηση; (σελ. 11)
      \item Τι σημαίνει ότι ένα στερεό σώμα εκτελεί σύνθετη κίνηση; (σελ. 12)
      \item Πότε ένα στερεό σώμα εκτελεί ομαλή περιστροφική κίνηση και ποια σχέση συνδέει το μέτρο της γωνιακής ταχύτητας με τη γωνία στροφής του σώματος στην περίπτωση αυτή; (σελ. 16)
      \item Πότε ένα στερεό σώμα εκτελεί ομαλά μεταβαλλόμενη περιστροφική κίνηση; (σελ. 16)
      \item Ποιες είναι οι χρονικές εξισώσεις που χρησιμοποιούνται στην ομαλά μεταβαλλόμενη περιστροφική κίνηση ενός στερεού σώματος; (σελ. 16-17)
      \item Με ποιον τρόπο μπορεί να υπολογιστεί γραφικά η γωνία στροφής (ή γωνιακή μετατόπιση) ενός στερεού σώματος; (σελ. 17)
      \item Τι ονομάζεται κέντρο μάζας ενός στερεού σώματος; (σελ. 17-18)
      \item Ένας τροχός εκτελεί σύνθετη κίνηση η οποία μπορεί να αναλυθεί σε δύο επιμέρους κινήσεις, μια μεταφορική και μια περιστροφική γύρω από άξονα που διέρχεται από το κέντρο μάζας του. Πώς μπορούμε να υπολογίσουμε την ταχύτητα ενός υλικού σημείου του τροχού; (σελ. 18)
      \item Για τροχό ακτίνας $R$ που κυλίεται χωρίς να ολισθαίνει, ποια σχέση συνδέει το μέτρο της ταχύτητας του κέντρου μάζας του και το μέτρο της γωνιακής του ταχύτητας; Ποια είναι η απόδειξη της σχέσης αυτής; (σελ. 18-19)
      \item Ένας τροχός κυλίεται χωρίς να ολισθαίνει σε οριζόντιο δάπεδο με ταχύτητα μέτρου $v_{\text{cm}}$. Να αποδείξετε ότι:
            \begin{itemize}
                  \item[α)] το ανώτατο σημείο της περιφέρειας του τροχού, ο οποίος βρίσκεται σε επαφή με το δάπεδο, έχει ταχύτητα μέτρου $2v_{\text{cm}}$,
                  \item[β)] το σημείο επαφής του τροχού με το δάπεδο έχει ταχύτητα ίση με μηδέν. (σελ. 19-20)
            \end{itemize}
      \item Ένας τροχός κυλίεται χωρίς να ολισθαίνει. Ποια σχέση συνδέει το μέτρο της επιτρόχιας (ή γραμμικής) επιτάχυνσης των σημείων της περιφέρειας του τροχού με το μέτρο της επιτάχυνσης του κέντρου μάζας του; Να αποδείξετε τη σχέση αυτή. (σελ. 20)
      \item Πώς μπορεί να γίνει η μελέτη της σύνθετης κίνησης ενός στερεού σώματος; (σελ. 21)
      \item Τι εκφράζει η ροπή μιας δύναμης; (σελ. 21)
      \item Πώς ορίζεται η ροπή δύναμης ως προς άξονα περιστροφής; (σελ. 21-22)
      \item Πώς ορίζεται η ροπή δύναμης ως προς σημείο; (σελ. 22)
      \item Πότε η ροπή μιας δύναμης ως προς κάποιον άξονα ή ως προς ένα σημείο ισούται με μηδέν; (σελ. 22)
      \item Τι κίνηση θα εκτελέσει ένα αρχικά ακίνητο ελεύθερο στερεό σώμα, αν ασκηθεί σε αυτό μια δύναμη; (σελ. 22-23)
      \item Τι ονομάζεται ζεύγος δυνάμεων; (σελ. 23)
      \item Πώς μπορούμε να υπολογίσουμε τη ροπή ενός ζεύγους δυνάμεων; (σελ. 23)
      \item Εξαρτάται η ροπή ενός ζεύγους δυνάμεων από το σημείο ή τον άξονα περιστροφής ως προς τον οποίο υπολογίζεται; Να αιτιολογήσετε την απάντησή σας. (σελ. 23)
      \item Αν σ' ένα αρχικά ακίνητο ελεύθερο στερεό σώμα ασκηθεί ένα ζεύγος δυνάμεων, τι κίνηση θα εκτελέσει το σώμα αυτό; (σελ. 24)
      \item Ποια είναι η συνθήκη ισορροπίας ενός στερεού σώματος που παραμένει ακίνητο; (σελ. 24)
      \item Πώς ορίζεται η στροφορμή ενός υλικού σημείου που εκτελεί κυκλική κίνηση; (σελ. 25)
      \item Πώς υπολογίζεται η στροφορμή ενός συστήματος υλικών σημείων; (σελ. 25)
      \item Ποια είναι η διατύπωση του θεμελιώδους νόμου της στροφικής κίνησης ενός υλικού σημείου; (σελ. 25)
      \item Με τι ισούται η αλγεβρική τιμή του ρυθμού μεταβολής της στροφορμής ενός συστήματος υλικών σημείων; (σελ. 25-26)
      \item Σε ποιες περιπτώσεις παραμένει σταθερή η στροφορμή ενός υλικού σημείου, το οποίο περιφέρεται γύρω από σταθερό άξονα; (σελ. 26)
      \item Σε ποιες περιπτώσεις παραμένει σταθερή η στροφορμή ενός συστήματος υλικών σημείων; (σελ. 26)
\end{enumerate}

\subsection{Κινηματική του στερεού σώματος -- συστήματος σωμάτων}

\begin{enumerate}
      \item Ποια είναι η σχέση που συνδέει τη γωνία περιστροφής $\Delta\theta$ ενός στερεού σώματος με το μέτρο της γωνιακής του ταχύτητας στην ομαλή περιστροφική κίνηση; (σελ. 56)
      \item Τι κίνηση εκτελούν όλα τα υλικά σημεία ενός στερεού σώματος (εκτός από αυτά του ακλόνητου άξονα), το οποίο εκτελεί ομαλή περιστροφική κίνηση; Ποια είναι η επιτάχυνση καθενός από τα σημεία αυτά και πώς υπολογίζεται το μέτρο της; (σελ. 56)
      \item Πώς υπολογίζουμε τον αριθμό των περιστροφών $N$ ενός στερεού σώματος όταν γνωρίζουμε τη γωνία στροφής $\Delta\theta$; (σελ. 58)
      \item Πώς μπορούμε να προσδιορίσουμε τη φορά της γωνιακής επιτάχυνσης ενός στερεού σώματος; Να εξηγήσετε αναλυτικά. (σελ. 58)
      \item Αν ένα στερεό σώμα εκτελεί μεταβαλλόμενη περιστροφική κίνηση, πώς υπολογίζουμε το μέτρο της επιτάχυνσης ενός υλικού σημείου του σώματος αυτού εκτός από αυτά του άξονα περιστροφής; (σελ. 61)
      \item Αν σχηματίσουμε τη γραφική παράσταση του μέτρου της γωνιακής ταχύτητας ενός στερεού σώματος σε συνάρτηση με το χρόνο, τότε τι παριστάνει το εμβαδόν που περικλείεται από τη γραφική παράσταση και τον άξονα του χρόνου από μια αρχική μέχρι μια τελική χρονική στιγμή; (σελ. 61)
      \item Σε ποιες επιμέρους κινήσεις μπορούμε να αναλύσουμε τη σύνθετη κίνηση που εκτελεί ένα στερεό σώμα; Τι σημαίνει ότι οι κινήσεις αυτές μπορεί να είναι ανεξάρτητες ή ''δεσμευμένες''; Τι μπορούμε να εφαρμόσουμε σε κάθε επιμέρους κίνηση; (σελ. 73)
      \item Αν ένα στερεό σώμα εκτελεί σύνθετη κίνηση, πώς μπορούμε να υπολογίσουμε την ταχύτητα ενός υλικού σημείου του; Να εξηγήσετε αναλυτικά. (σελ. 73)
      \item Ποιες είναι οι σχέσεις που συνδέουν τα μεταφορικά με τα περιστροφικά μεγέθη στην περίπτωση της κύλισης χωρίς ολίσθηση; Ποια είναι η ακτίνα $R$ που χρησιμοποιείται στις σχέσεις αυτές; Να εξηγήσετε αναφέροντας ένα παράδειγμα. (σελ. 74-75)
      \item Πώς μπορούμε να μελετήσουμε την περίπτωση της κύλισης ενός σώματος; (σελ. 75)
      \item Αν ένα στερεό σώμα εκτελεί κύλιση, τότε για ποια υλικά σημεία του η γραμμική ταχύτητά τους ισούται κατά μέτρο με την ταχύτητα του κέντρου μάζας; (σελ. 79)
      \item Αν ένα στερεό σώμα εκτελεί σύνθετη κίνηση, τότε πώς μπορούμε να υπολογίσουμε την επιτάχυνση ενός υλικού σημείου του; Να εξηγήσετε αναλυτικά. (σελ. 79)
      \item Πότε ισχύουν οι τύποι $s_{\text{cm}} = r \cdot \Delta\theta$, $v_{\text{cm}} = \omega r$ και $a_{\text{cm}} = a_{\text{γων}} r$ στην περίπτωση του γιο-γιο; (σελ. 83)
      \item Πού οφείλεται το τύλιγμα ή το ξετύλιγμα του νήματος στην περίπτωση ενός στερεού σώματος που εκτελεί σύνθετη κίνηση; Πώς υπολογίζεται το μήκος του νήματος που τυλίγεται ή ξετυλίγεται; (σελ. 83)
      \item Ποιο είναι το σύστημα 1ου είδους και πώς μπορούμε να μελετήσουμε την κίνησή του; (σελ. 96)
      \item Ποιο είναι το σύστημα 2ου είδους και πώς μπορούμε να μελετήσουμε την κίνησή του; (σελ. 99)
      \item Πώς βρίσκουμε τους κινηματικούς συνδέσμους της ταχύτητας στην περίπτωση της τροχαλίας; (σελ. 102)
      \item Πώς βρίσκουμε τους κινηματικούς συνδέσμους της ταχύτητας σε ένα σύστημα 2ου είδους; (σελ. 106)
      \item Ποιο είναι το σύστημα 3ου είδους και πώς μπορούμε να μελετήσουμε την κίνησή του; (σελ. 109)
\end{enumerate}

\subsection{Ροπή δύναμης -- ισορροπία του στερεού σώματος}

\begin{enumerate}
      \item Πώς βρίσκουμε το μοχλοβραχίονα μιας δύναμης; (σελ. 146)
      \item Πότε η ροπή μιας δύναμης ως προς άξονα ισούται με μηδέν; Πότε συμβαίνει το ίδιο στην περίπτωση της ροπής μιας δύναμης ως προς ένα σημείο; (σελ. 146)
      \item Πώς μπορούμε να υπολογίσουμε τη ροπή μιας δύναμης αν τη δύναμη αυτή την αναλύσουμε σε δύο κάθετες μεταξύ τους συνιστώσες; (σελ. 146)
      \item Τι ονομάζεται ζεύγος δυνάμεων; (σελ. 147)
      \item Πώς υπολογίζεται το μέτρο της ροπής ενός ζεύγους δυνάμεων; Εξαρτάται ο τύπος που γράψατε από τον άξονα ή το σημείο ως προς το οποίο υπολογίζεται η ροπή του ζεύγους; (σελ. 147)
      \item Όταν ένα στερεό σώμα είναι ακίνητο, ποιες σχέσεις ισχύουν για τις δυνάμεεις που δέχεται το σώμα και ποιες για τις ροπές των δυνάμεων αυτών; (σελ. 152)
      \item Πότε είναι χρήσιμο να αναλύουμε τις δυνάμεις που ασκούνται σ' ένα ακίνητο στερεό σώμα; Μας δεσμεύει η ανάλυση αυτή για την εφαρμογή της συνθήκης που ισχύει για τις ροπές των δυνάμεων αυτών; (σελ. 153)
      \item Πώς σχεδιάζουμε τη δύναμη που δέχεται ένα στερεό σώμα από: i) ένα λείο στήριγμα; ii) ένα λείο επίπεδο; iii) μια άρθρωση; iv) ένα αβαρές νήμα; (σελ. 154-155)
      \item Ως προς ποιο σημείο είναι προτιμότερο να υπολογίζουμε τη συνισταμένη των ροπών όταν χρησιμοποιούμε τη σχέση $\Sigma\tau = 0$; Για ποιο λόγο; (σελ. 157)
      \item Ποια είναι τα συχνότερα λάθη που γίνονται στην εφαρμογή της σχέσης $\Sigma\tau = 0$; Πώς μπορούμε να τα αποφύγουμε; (σελ. 157)
      \item Από τι εξαρτώνται η φορά και το μέτρο της δύναμης που δέχεται από την άρθρωση ένα στερεό σώμα όταν ισορροπεί; Τι πρέπει να σχεδιάζουμε πρώτα για το λόγο αυτό; (σελ. 157)
      \item Σε τι αντιστέκεται η στατική τριβή που ασκείται μεταξύ δύο επιφανειών; Τι τιμές μπορεί να πάρει το μέτρο της στατικής τριβής; (σελ. 160)
      \item Τι συμβαίνει με τις δυνάμεις από τα δάπεδα, τους τοίχους, τα νήματα και τις αρθρώσεις όταν αλλάξει κάποια δύναμη που δέχεται το στερεό σώμα το οποίο ισορροπεί; (σελ. 160)
      \item Πώς μπορούμε να μελετήσουμε την ισορροπία ενός συστήματος σωμάτων που απαρτίζεται από δύο ή περισσότερα σώματα τα οποία είναι μεταξύ τους κολλημένα; (σελ. 165)
      \item Αν διαθέτουμε ένα σύστημα σωμάτων που ισορροπεί και το οποίο απαρτίζεται από σώματα που είναι κολλημένα μεταξύ τους (σύστημα 1ου είδους), τι πρέπει να κάνουμε αν θέλουμε να βρούμε κάποια εσωτερική δύναμη του συστήματος; (σελ. 165)
      \item Ποιοι είναι οι δύο τρόποι που μπορούμε να χρησιμοποιήσουμε για να μελετήσουμε την ισορροπία ενός συστήματος σωμάτων τα οποία δεν είναι κολλημένα μεταξύ τους; (σελ. 168)
      \item Όταν μια ράβδος ή μια σανίδα ισορροπεί οριζόντια με τη βοήθεια δύο υποστηριγμάτων, τι πρέπει να ελέγξουμε ώστε να αποκλείσουμε την ανατροπή της ράβδου ή της σανίδας; (σελ. 170-171)
      \item Πώς πρέπει να μελετήσουμε την περίπτωση όπου έχουμε ένα σύστημα σωμάτων για το οποίο κάποιο ή κάποια σώματα κινούνται, ενώ τα υπόλοιπα παραμένουν ακίνητα; (σελ. 181)
\end{enumerate}

\subsection{Στροφορμή και τη διατήρηση της στροφορμής}

\begin{enumerate}
      \item Με τι ισούται ο ρυθμός μεταβολής της ορμής ενός υλικού σημείου και με τι ο ρυθμός μεταβολής της στροφορμής του; Τι πρέπει να αναφέρουμε όταν υπολογίζουμε το ρυθμό μεταβολής της στροφορμής; (σελ. 230)
      \item Πώς μπορούμε να βρούμε τη μεταβολή της στροφορμής ενός υλικού σημείου; Τι μπορούμε να κάνουμε όταν η αρχική και η τελική στροφορμή του υλικού σημείου βρίσκονται στην ίδια διεύθυνση; (σελ. 230)
      \item Πώς υπολογίζεται η στροφορμή ενός συστήματος υλικών σημείων (ή γενικότερα σωμάτων); (σελ. 235)
      \item Πότε μπορούμε να γράψουμε για ένα υλικό σημείο ότι $\frac{d\vec{L}}{dt} = \frac{\Delta\vec{L}}{\Delta t}$; (σελ. 235)
\end{enumerate}

\end{document}